% answers_short.tex —— 2003 年·粤桂卷 · 简短答案速查
% 由 tools/make_short_answers.py 自 answers.tex 的 \gksolhead/\gkanswer 提取，
% 并与 2003/2003_数学真题_粤桂_答案要点.md 逐题比对后生成，勿手改（改源 MD 后重跑）。

\begin{gknotice}
\textbf{使用说明}：本节只给最终答案便于快速核对，完整推导见解析版。答案为编写者独立推导，\textbf{非官方标准答案}。
\end{gknotice}

\begin{gknotice}
选择题答案中的 U 表示待核，未给出猜测答案，详情见解析版。
\end{gknotice}

\gkpart{选择题}
\begin{center}\begin{tabular}{ccccccccccccc}
\toprule
题号 & 1 & 2 & 3 & 4 & 5 & 6 & 7 & 8 & 9 & 10 & 11 & 12 \\
\midrule
答案 & A & D & C & C & B & D & A & C & B & U & C & A \\
\bottomrule
\end{tabular}\end{center}

\gkpart{填空题}
\begin{itemize}[leftmargin=2.6em]
  \item[\textbf{13.}] $(2, 4]$
  \item[\textbf{14.}] $-\dfrac{21}{2}$
  \item[\textbf{15.}] $S_{\triangle ABC}^2 + S_{\triangle ACD}^2 + S_{\triangle ADB}^2 = S_{\triangle BCD}^2$
  \item[\textbf{16.}] 72
\end{itemize}

\gkpart{解答题}
\begin{itemize}[leftmargin=2.6em]
  \item[\textbf{17.}] (1) 证明见解析；(2) $\dfrac{2\sqrt{3}}{3}$
  \item[\textbf{18.}] $|z| = \sqrt{2} - 1$
  \item[\textbf{19.}] $\left(0, \dfrac{1}{2}\right] \cup [1, +\infty)$
  \item[\textbf{20.}] 12 小时后
  \item[\textbf{21.}] 当 $0 < a < \dfrac{\sqrt{2}}{2}$ 时，存在两定点 $\left(\pm\sqrt{\dfrac{1}{2} - a^2}, a\right)$，距离之和为定值 $\sqrt{2}$；当 $a = \dfrac{\sqrt{2}}{2}$ 时，点 $P$ 轨迹为圆，两定点重合于圆心 $\left(0, \dfrac{\sqrt{2}}{2}\right)$，距离之和为 $\sqrt{2}$；当 $a > \dfrac{\sqrt{2}}{2}$ 时，存在两定点 $\left(0, a \pm\sqrt{a^2 - \dfrac{1}{2}}\right)$，距离之和为定值 $2a$。
  \item[\textbf{22.}] (1) 证明见解析；(2) $0 < a_0 < \dfrac{1}{3}$
\end{itemize}
